% Pursuing Stacks (cotes 134-2 to 134-8): archive cote 134-6, site page 17 (PDF page 18), typescript page 336.
% typescript. Draft compared: ps5.tex:2221-2291.
% AI-transcribed by Claude Opus 5.5 (claude-opus-5-5, adaptive thinking, effort high), Message Batches API, batch ps3 (msgbatch_01XdFmuVVBAcgc7EFxdtCqsj), 2026-11-12; not yet checked by eye.
\scanpage{134-6/17}
\mspage{336}

(22) \[ L_\bullet^{\hand{\omega}}\deleted{\unsure{(\omega)}} = ((L_n)_{n\in\underline{N}} \ , (d_n^{\hand{\omega}})_{n\in\underline{N}} ) \]
associated to \xout{the} all possible systems of orientations
(23) \[ \hand{\omega} = (\hand{\omega}_a)_{a\, \hand{\in}\, \hand{O}b(A)} \]
of the various cells  i(a)  ,  there is a canonical \uline{transitive system}
of isomorphisms,  by defining the isomorphism
(24) \[ u_{\hand{\omega,\omega'}} : L_\bullet \hand{\xrightarrow{\sim}} L_\bullet' \]
for two different choices \hand{$\omega, \omega'$} of systems of orientations,  by
\margin{\hand{$\varepsilon_a^{\omega,\omega'}$} (circled)}
(25) \[ u_{\hand{\omega,\omega'}} \hand{\mid} a = \hand{\varepsilon}_a^{\hand{\omega,\omega'}} \mathrm{id}_a \ , \qquad \hand{\varepsilon}_a^{\hand{\omega,\omega'}} \hand{\in\{} \pm 1\hand{\}}\hand{,} \]
where  \uline{the sign}  \hand{$\varepsilon$}$_a^{\hand{\omega,\omega'}}$ is  \uline{equal to}\hand{+}1 \uline{or} -1,  \uline{according to whether}  \hand{$\omega$}$_a$  \uline{and}
\hand{$\omega$}$'_a$  \uline{are equal or not}.  It is immediate that  (24)  then is an isomor-%
phism\xout{s} componentwise,  respecting degrees,  and commu\xout{ta}ting to the res-%
pective differential  operators.  Transitivity of the isomorphisms
(24) for a triple  (\hand{$\omega$} ,\hand{$\omega'$},\hand{$\omega''$})  is equally immediate.  This implies that
by this transitive system of isomorphisms,  \uline{we may identify all the chai\hand{n}}
\uline{complexes}  $L_\bullet^{\hand{\omega}}$  in Add(A$^{o}$) to a single chain complex,  canonically \deleted{\xout{equiv}}
isomorphic to each $L_\bullet^{\hand{\omega}}$ ,  and which we may just designate by L. .  This
chain complex now is defined intrinsically in terms of the category
structure of A (up to canonical isomorphism),  in the ``safe'' case at
any rate when every isomorphism in A is an identity,  so that in the
construction of $L_\bullet^{\hand{\omega}}$ there enters no other choice besides \hand{$\omega$}.  Otherwise\unsure{,}
as seen above (precedent page),  we must still suitably choose the so-%
called ``strict'' boundary operators  (13) ,  among all monomorphisms
b \hand{$\hookrightarrow$} a  in A such that  dim a = dim b + 1 .

In the first case say (isomorphisms being identities),  all con-%
ditions considered for A are stable under finite products,  that's why
in terms of the three standard cases of \hand{$\Delta$}, \hand{$\square$} and \hand{$\bigcirc$},  we could con-%
struct others by taking finite products.  The three standard test catego-%
ries \xout{xxx} may be viewed as particularly ``economic'' or skillful ways of
``cutting out'' a suitable bunch of cellular decompositions\hand{,} and \inserted{of} elimina-%
ting automorphisms (by total ordering of vertices and the like ...),
so as to insure\hand{:} a) that isomorphisms in A are identities ,  b)  the
canonical chain complex L.  in Add(A$^{o}$)  is an abelianizator,  and  c)
A moreover is a strict test category,  and even a contractor.  On the
other hand,  as all these conditions  (plus the condition (15) of course
about the i(a)' s  representing n-cells)  \deleted{all these \xout{conditions are stab}}

\margin{17}
